% !PRINTED 637
\weberpage{637}
Diese Formeln gelten für jedes $\nu$, wenn wir auch noch in Bezug auf die
$F_\nu$ festsetzen, dass $F_h = F_k$ sein soll, wenn $h \equiv k
\pmod{n-1}$ ist. Es ergiebt sich dann aus 4., dass die Functionen $F_0,
F_1 \ldots F_{n-2}$ durch die Substitution $s$ ungeändert bleiben, und
durch $t^{-1}$ eine cyklische Permutation erfahren.

Wenn wir hier die Exponenten auf ihre kleinsten positiven Reste nach dem
Modul $n$ reduciren wollen, so setzen wir

\weq{15}{g^\nu = n q_\nu + r_\nu, \quad 0 < r_\nu < n, \quad r_0 = 1,}

und die Zahlen $r_\nu$ fallen in irgend einer Reihenfolge mit den Zahlen
$1, 2 \ldots n-1$ zusammen, so dass $r_0$ immer $= 1$ ist.

Dann wird nach (5)

\weq{16}{f_{\nu-1} = (\varepsilon^{r_\nu}, x)(\varepsilon^{r_{\nu-1}}, x)^{-g},}

und wir setzen noch

\weq{17}{\Phi_\nu = F_\nu f_\nu^{q_{n-2}} f_{\nu+1}^{q_{n-3}} \cdots f_{\nu+n-3},}

und erhalten aus (14)

\weq{18}{(\varepsilon^{r_\nu}, x)^n = \Phi_\nu^n f_\nu^{r_{n-2}} f_{\nu+1}^{r_{n-3}} \cdots f_{\nu+n-2}^{r_0}.}

Die Exponenten $r_\nu$ bleiben ungeändert, wenn $g$ durch $g - l n$
ersetzt wird, sind also von der Bedingung (10) unabhängig und \gesperrt{können aus
einer beliebigen primitiven Wurzel $g$ abgeleitet} werden.

Die charakteristischen Eigenschaften der in diesen Formeln vorkommenden
Functionen $f_\nu, F_\nu, \Phi_\nu$, deren Gesammtheit wir mit $\omega_\nu$
bezeichnen wollen, sind, um es nochmals zu wiederholen, folgende:

$\alpha)$ \gesperrt{Durch die cyklische Permutation $s = (h, h+1)$ unter den
Indices von $x$ ändern sich die Functionen $\omega_\nu$ nicht.}

$\beta)$ \gesperrt{Durch die lineare Permutation $t^{-1} = (h, g^{-1} h)$ unter den
Indices der $x$ wird unter den Indices der $\omega$ die cyklische
Permutation $(0, 1, 2 \ldots n-2)$ hervorgerufen.}

$\gamma)$ \gesperrt{Durch die Substitution $\sigma = (\varepsilon, \varepsilon^g)$
wird unter den Indices von $\omega$ dieselbe cyklische Permutation
$(0, 1, 2 \ldots n-2)$ bewirkt.}

$\delta)$ \gesperrt{Die cyklischen Functionen von $\omega_\nu$ sind metacyklische
Functionen der $x$ und von $\varepsilon$ frei.}
